% =====================================================================
%  Information and Coding 2026/27 — Lab work n.º 1
%  Relatório 2: utilização de ferramentas de IA (máximo 6 páginas)
% =====================================================================
\documentclass[a4paper,11pt]{article}

\usepackage[utf8]{inputenc}
\usepackage[T1]{fontenc}
\usepackage[portuguese]{babel}
\usepackage[margin=2.2cm]{geometry}
\usepackage{booktabs}
\usepackage{tabularx}
\usepackage{xcolor}
\usepackage{hyperref}
\hypersetup{colorlinks=true,linkcolor=blue!50!black,urlcolor=blue!50!black}

\newcommand{\todo}[1]{\textcolor{red}{\textbf{[TODO:} #1\textbf{]}}}
\newcommand{\prog}[1]{\texttt{#1}}

% ---------------------------------------------------------------------
% Modelo de caso de uso. Responde às perguntas do enunciado:
% problema, proposta da IA, verificação, correção, alterações, aprendizagem.
% Uso: \casoIA{Título}{Ferramenta}{Membro}{Problema}{Proposta}{Verificação}
%              {Correto?}{Alterações}{Aprendizagem}
% ---------------------------------------------------------------------
\newcommand{\casoIA}[9]{%
  \subsubsection*{#1}
  \noindent\textit{Ferramenta:} #2 \quad \textit{Membro:} #3
  \begin{description}\itemsep2pt
    \item[Problema.] #4
    \item[Proposta da IA.] #5
    \item[Verificação.] #6
    \item[Estava correto?] #7
    \item[O que alterámos.] #8
    \item[O que aprendemos.] #9
  \end{description}}

\title{Information and Coding --- Trabalho Prático n.º 1\\[4pt]
\large Relatório 2: Utilização de ferramentas de Inteligência Artificial}
\author{%
  Nome Apelido 1 (n.º mec.~XXXXX) \and
  Nome Apelido 2 (n.º mec.~XXXXX) \and
  Nome Apelido 3 (n.º mec.~XXXXX)}
\date{Universidade de Aveiro --- DETI\\ Outubro de 2026}

\begin{document}
\maketitle

% =====================================================================
\section{Introdução}
% =====================================================================
\subsection{Objetivo do relatório}
\todo{Descrever o propósito e o âmbito.}

\subsection{Metodologia de registo}
\todo{Como foram registadas as interações ao longo do trabalho
(p.\,ex.\ ficheiro partilhado no repositório, data, membro, ferramenta, tarefa).}

% =====================================================================
\section{Ferramentas de IA utilizadas}
% =====================================================================
\begin{table}[h]
\centering
\caption{Ferramentas de IA utilizadas e finalidade.}
\label{tab:ferramentas}
\begin{tabularx}{\linewidth}{llX}
\toprule
Ferramenta & Membro(s) & Utilização principal \\
\midrule
Claude (Anthropic) & -- & \todo{p.\,ex.\ planeamento, estrutura do relatório} \\
GitHub Copilot & -- & \todo{p.\,ex.\ autocompletar código} \\
ChatGPT & -- & \todo{} \\
\bottomrule
\end{tabularx}
\end{table}

% =====================================================================
\section{Casos de utilização}
% =====================================================================
% Escolher os casos mais relevantes (sucessos E falhas) — 6 páginas é pouco.

\subsection{Planeamento e escolhas tecnológicas}
\casoIA{Escolha das linguagens de programação}
  {Claude}{\todo{nome}}
  {Decidir que linguagens usar em cada parte do trabalho.}
  {\todo{p.\,ex.\ C++ para os programas e codecs, Python para gráficos e benchmarks.}}
  {\todo{}}{\todo{}}{\todo{}}{\todo{}}

\subsection{Parte I --- Ferramentas de áudio}
\casoIA{\todo{Título do caso}}
  {\todo{}}{\todo{}}{\todo{}}{\todo{}}{\todo{}}{\todo{}}{\todo{}}{\todo{}}

\subsection{Parte II --- Codec sem perdas}
\casoIA{\todo{p.\,ex.\ Implementação do código de Golomb}}
  {\todo{}}{\todo{}}{\todo{}}{\todo{}}{\todo{}}{\todo{}}{\todo{}}{\todo{}}

\subsection{Parte III --- Codec com perdas}
\casoIA{\todo{p.\,ex.\ Quantização dos coeficientes DCT}}
  {\todo{}}{\todo{}}{\todo{}}{\todo{}}{\todo{}}{\todo{}}{\todo{}}{\todo{}}

\subsection{Depuração, testes e medição de desempenho}
\casoIA{\todo{Título do caso}}
  {\todo{}}{\todo{}}{\todo{}}{\todo{}}{\todo{}}{\todo{}}{\todo{}}{\todo{}}

\subsection{Escrita dos relatórios}
\casoIA{\todo{p.\,ex.\ Estrutura em \LaTeX{} dos relatórios}}
  {Claude}{\todo{}}{\todo{}}{\todo{}}{\todo{}}{\todo{}}{\todo{}}{\todo{}}

% =====================================================================
\section{Avaliação crítica}
% =====================================================================
\subsection{Onde a IA foi mais útil}
\todo{}

\subsection{Erros e limitações observados}
\todo{Código incorreto, alucinações, soluções ineficientes, detalhes subtis
(p.\,ex.\ divisão inteira com números negativos, formato WAV).}

\subsection{Estratégias de verificação}
\todo{Testes de ida e volta (\prog{cmp} entre o original e o descodificado),
comparação com ferramentas de referência, revisão manual.}

\subsection{Impacto na produtividade e na aprendizagem}
\todo{Tempo poupado vs.\ tempo gasto a corrigir; o que teríamos feito de outra forma.}

% =====================================================================
\section{Conclusões}
% =====================================================================
\todo{Balanço global e recomendações para uso futuro.}

\end{document}
